%!TEX program = lualatex
\documentclass[11pt]{article}
\usepackage{hyperref}
\usepackage{hyphenat}
\usepackage{fontspec}
\setmainfont{Carlito}
\usepackage[scaled=.5]{helvet}
\usepackage{courier}
\usepackage{amsmath}
\usepackage{booktabs}
\usepackage{setspace}
\usepackage{geometry}
\usepackage{fancyvrb}
\setstretch{0.97}
\geometry{margin=1in}

\title{\textbf{Enterprise Risk Management --- Exercises \& Case Studies}}
\author{}
\date{}

\begin{document}
\maketitle

A mid-sized bank operates a trading desk managing a diversified market-risk portfolio with a total market value of \textbf{€100 million}. The bank has implemented an ERM framework centered on quantitative risk limits, with Value-at-Risk (VaR) as the primary risk metric.

\bigskip
\hrule
\bigskip

\section*{How to use this student workbook}
This student workbook contains a merged, pedagogical set of exercises on risk measurement (VaR and ES) and applied ERM case studies. Each exercise begins with a short contextual introduction that explains the learning objective. Where appropriate, you will be asked to work with a realistic, larger dataset; instructions for generating and analysing the data in \textbf{Excel} are provided. No solutions are included in this file.

\newpage
\part*{Part I --- Risk Measurement Foundations}

\section{Lecture 1: Intuition about VaR and ES}
\subsection*{Learning objective}
Understand what VaR and Expected Shortfall (ES) measure, their limitations, and how sample size and tail behavior affect them.

\subsection*{Exercise 1.1 --- Conceptual VaR}
A trading desk reports a \textbf{1-day VaR at 99\% confidence} of \textbf{€10 million}.

\begin{enumerate}
  \item In plain language, explain what this number means.
  \item Over 100 trading days, how many days would you expect losses to exceed €10 million?
  \item Does VaR tell you how large losses can be beyond €10 million? Explain.
  \item Give two short examples of situations where VaR can be misleading.
\end{enumerate}

\bigskip
\subsection*{Exercise 1.2 --- Historical VaR with a realistic sample (Excel)}
\textbf{Context:} Small samples (10 observations) are useful for demonstrating mechanics but not for building intuition about model stability. In this exercise you will create a realistic 250-day return series in Excel with occasional large negative shocks and compute historical VaR and ES.

\textbf{Excel data generation (no code required):}
\begin{enumerate}
  \item Open a new Excel workbook. In column A (A2:A251) create 250 rows labelled Day 1 to Day 250.
  \item In column B (B2:B251) generate daily returns using Excel's random functions and a simple volatility clustering approximation:
    \begin{itemize}
      \item In B2 enter: \texttt{=NORM.INV(RAND(),0.0004,0.01)} to create a base normal return.
      \item In C2 enter a volatility factor: \texttt{=1 + 0.5*ABS(NORM.INV(RAND(),0,0.8))}.
      \item In D2 compute the adjusted return: \texttt{=B2*C2}.
      \item Copy B2:D2 down to row 251.
    \end{itemize}
  \item To inject occasional negative jumps, choose five row indices (for example rows 32, 80, 122, 192, 232) and in column D for those rows subtract a jump value (e.g., -0.05, -0.08, or -0.12). Example: if D32 currently contains a value, replace it with \texttt{=D32 - 0.08}.
  \item Save the sheet as \texttt{returns.xlsx}.
\end{enumerate}

\textbf{Excel calculations for VaR and ES:}
\begin{enumerate}
  \item In column E compute losses as negative returns: \texttt{= -D2} and copy down.
  \item Sort a copy of column E in descending order (largest losses first) using Excel's \texttt{SORT} or by copying to a new sheet and using Data → Sort.
  \item 95\% historical VaR: locate the 95th percentile of losses using \texttt{=PERCENTILE.EXC(E2:E251,0.95)}.
  \item 99\% historical VaR: \texttt{=PERCENTILE.EXC(E2:E251,0.99)}.
  \item 95\% ES: compute the average of losses greater than or equal to the 95\% VaR using \texttt{=AVERAGEIF(E2:E251,">=" & cell_with_VaR95)}.
  \item 99\% ES: \texttt{=AVERAGEIF(E2:E251,">=" & cell_with_VaR99)}.
  \item Convert to euros for a €50m portfolio: multiply the loss fractions by 50,000,000.
\end{enumerate}

\textbf{Deliverables:} a short report (1 page) with the computed VaR/ES numbers (in \euro), a histogram (Excel chart of column D), the backtest count (number of days where loss > VaR99), and a 3–4 sentence interpretation.

\newpage
\section{Lecture 2: Model risk and stress testing}
\subsection*{Learning objective}
Learn how model assumptions (normality, stationarity) and stress scenarios affect risk metrics and decision-making.

\subsection*{Exercise 2.1 --- Parametric VaR vs Historical VaR (Excel)}
\begin{enumerate}
  \item Using the synthetic dataset in \texttt{returns.xlsx}, compute the sample standard deviation of returns with \texttt{=STDEV.S(D2:D251)}.
  \item Parametric 99\% VaR (assuming normality) = \texttt{=NORM.S.INV(0.99) * stdev}, where \texttt{stdev} is the cell with the standard deviation.
  \item Compare the parametric VaR to the historical VaR computed earlier and write a short explanation of the difference.
  \item Stress test: compute the standard deviation of the most recent 20 days with \texttt{=STDEV.S(D232:D251)} and double it. Recompute parametric VaR using the doubled volatility and comment on the change.
\end{enumerate}

\subsection*{Exercise 2.2 --- Scenario analysis (Excel)}
\begin{enumerate}
  \item Construct a stress scenario for the €100m portfolio in Case 1: specify an equity shock (e.g., -8\%), an interest-rate shock (e.g., +30 bps), and a liquidity shock (e.g., margin calls equal to X\% of buffer). Use Excel to compute the P\&L impact:
    \begin{itemize}
      \item Equity P\&L = 60,000,000 * equity shock.
      \item Interest-rate P\&L = sensitivity per bp * bp shock.
      \item Liquidity impact: specify a margin outflow and compute required asset sales if buffer is exceeded.
    \end{itemize}
  \item Explain in a short paragraph how scenario analysis complements VaR and ES in ERM.
\end{enumerate}

\newpage
\part*{Part II --- Applied ERM Case Studies}

\section{Case 1: Market risk, VaR limits and stress}
\subsection*{Context (excerpt)}
A mid-sized bank operates a trading desk managing a diversified market-risk portfolio with a total market value of \textbf{€100 million}. The bank has implemented an ERM framework centered on quantitative risk limits, with Value-at-Risk (VaR) as the primary risk metric.

\subsection*{Portfolio composition and parameters}
\begin{itemize}
  \item Equity exposure: \euro 60m; daily equity volatility: 2\%.
  \item Interest-rate exposure: \euro 40m; daily yield volatility: 5 bps; sensitivity: +1 bp $\Rightarrow$ \euro 0.20m loss.
  \item Correlation between equity returns and yield changes: 0.
  \item 1-day 99\% VaR limit: \euro 3.0m.
\end{itemize}

\subsection*{Exercise 3.1}
\begin{enumerate}
  \item Compute the 1-day 99\% VaR of the portfolio using a variance--covariance approach (show intermediate steps in Excel).
  \item Does the desk comply with the VaR limit? Explain.
  \item Consider a stress scenario: equity decline 8\%, yields +30 bps. Compute the scenario loss in Excel.
  \item Explain why such a stress loss can occur even if VaR is within limits.
  \item Propose two complementary ERM controls (non-VaR) that would reduce the chance of surprise losses.
\end{enumerate}

\bigskip
\section{Case 2: Expected Shortfall, governance and escalation}
\subsection*{Context}
The bank now monitors 1-day 99\% Expected Shortfall (ES) with a limit of \euro 4.0m and mandatory escalation on breach. The last 12 trading days' P\&L (in \euro m) are:

\begin{center}
\begin{tabular}{cccc}
\toprule
Day & P\&L (\euro m) & Day & P\&L (\euro m) \\
\midrule
1 & +0.6 & 7 & -3.5 \\
2 & -1.1 & 8 & +0.5 \\
3 & +0.4 & 9 & -1.7 \\
4 & -2.3 & 10 & -4.2 \\
5 & +0.9 & 11 & -0.6 \\
6 & -0.8 & 12 & -6.8 \\
\bottomrule
\end{tabular}
\end{center}

\subsection*{Exercise 3.2}
\begin{enumerate}
  \item Sort the P\&L distribution from worst to best (Excel: use SORT or sort the column).
  \item Compute the empirical 1-day 99\% ES and determine whether the ES limit is breached (Excel: use \texttt{=AVERAGEIF(range,">=" & VaR99)} after computing VaR99 with \texttt{PERCENTILE.EXC}).
  \item The large loss on Day 12 occurred without prior escalation. Identify possible failures in: risk reporting, governance architecture, and incentives/risk culture.
  \item Explain why replacing VaR with ES alone does not guarantee effective ERM.
  \item Propose a practical escalation rule (quantitative + qualitative) that would have forced earlier action in this scenario.
\end{enumerate}

\newpage
\section{Case 3: Signals without escalation (governance and culture)}
\subsection*{Context}
A bank expands rapidly into structured finance. Risk reports show rising concentration, frequent exceptions, and stress-test sensitivity. Committee minutes use soft language and leave follow-up to business judgment. A downturn causes large losses and supervisory scrutiny.

\subsection*{Exercise 3.3}
\begin{enumerate}
  \item Based on the description, identify the type of ERM framework (principles-based vs rules-based) and justify your answer.
  \item Reconstruct the ERM architecture: who takes risk, who monitors, who decides escalation?
  \item List the early warning signals and explain why they failed to trigger corrective action.
  \item Explain how incentives and culture can undermine ERM, giving two concrete examples.
  \item Propose two concrete governance changes that would make escalation automatic when certain signals appear.
\end{enumerate}

\newpage
\section{Case 4: Hedging effectiveness (prospective vs retrospective)}
\subsection*{Context}
A firm hedges commodity exposure using futures. Normal-period estimates (24 months) and stress-period realized statistics are provided. The firm uses the minimum-variance hedge ratio and evaluates hedging effectiveness by variance reduction. Margining creates liquidity exposure.

\subsection*{Data (monthly, \euro m)}
\begin{center}
\begin{tabular}{lcc}
\toprule
 & Normal period & Stress period \\
\midrule
E[X] & 0.0 & -2.0 \\
Var(X) & 100 & 225 \\
Var(H) & 64 & 144 \\
Cov(X,H) & 48 & 12 \\
\bottomrule
\end{tabular}
\end{center}

Initial margin: \euro 10m; average monthly variation margin outflow: \euro 18m; liquidity buffer: \euro 25m.

\subsection*{Exercise 3.4}
\begin{enumerate}
  \item Compute the minimum-variance hedge ratio $h^*$ using normal-period statistics (Excel: compute \texttt{=Cov/VarH}).
  \item Using $h^*$, compute prospective hedged variance and hedging effectiveness (HE$_{pros}$) in Excel.
  \item Using the same $h^*$ but stress-period statistics, compute retrospective hedged variance and HE$_{retro}$.
  \item Explain why HE$_{retro}$ differs from HE$_{pros}$ in terms of correlation, basis risk, and regime change.
  \item Is the liquidity buffer sufficient on average? What governance changes would reduce the probability of forced asset sales if margin spikes to \euro 40m?
  \item Identify three risk categories that interacted in this episode and explain how ERM should govern their interaction.
\end{enumerate}

\newpage
\section{Case 5: When hedging increased risk}
\subsection*{Context}
A large industrial firm hedges commodity price risk with derivatives. Ex ante the hedge reduces variance by 75\% (h* = 0.9). After a shock, correlations weaken, volatility spikes, and margin calls create liquidity stress. HE$_{retro} < 0$.

\subsection*{Exercise 3.5}
\begin{enumerate}
  \item Identify the main risks and how they interact.
  \item Explain why the hedge looked effective ex ante.
  \item Interpret the hedge ratio $h^*$ and its assumptions.
  \item Explain why prospective effectiveness failed to predict outcomes.
  \item Which secondary risks were insufficiently considered?
  \item Propose two ERM governance changes that would have reduced the crisis severity.
\end{enumerate}

\newpage
\section{Case 6: Shell --- ERM at a large multinational}
\subsection*{Context (excerpt)}
Shell operates across upstream, integrated gas, chemicals, trading, power, and retail. Its ERM is principles-based, scenario-driven, and integrates risk into strategy and capital allocation.

\subsection*{Exercise 3.6}
Answer the following (short essay style, 200--400 words each):
\begin{enumerate}
  \item Characterize Shell's ERM framework: rules-based or principles-based? Compliance-oriented or decision-oriented? Support your answer with evidence.
  \item Why does Shell use scenario analysis for long-term risks? What are the limits of quantitative measurement?
  \item How is risk ownership and oversight allocated? When could this separation weaken risk management?
  \item Identify three principal risks and explain how Shell integrates them into strategy.
  \item Describe a plausible failure mode for Shell's ERM and the early warning signals you would monitor.
\end{enumerate}

\newpage
\part*{Part III --- Group ERM Design Exercise}

\section*{Group assignment (included)}
\subsection*{Learning objective}
Design a realistic ERM system for a complex, capital-intensive multinational. The exercise tests synthesis: governance, risk identification, measurement, response, and integration into strategy.

\subsection*{Task}
Prepare a 2--3 page ERM design proposal covering:
\begin{enumerate}
  \item ERM philosophy and scope (principles vs rules; focused vs comprehensive).
  \item ERM architecture and governance (who owns risks, who oversees, escalation).
  \item Risk identification and assessment (how to capture emerging and interacting risks).
  \item Risk response and instruments (how to govern derivatives, insurance, contingent liquidity).
  \item Integration into strategy and capital allocation.
  \item Reporting, dashboards, and escalation rules.
  \item Two plausible failure modes and early warning signals.
\end{enumerate}

\subsection*{Deliverable}
A 2--3 page written proposal and a 10-minute group presentation. Include a one-page dashboard mock-up (PDF or slide) showing key indicators and escalation triggers.

\bigskip
\hrule
\bigskip
\noindent\textbf{End of student workbook.}
\end{document}